\chapter{BCD译码器——5大变式详解}

\section{变式1：多位BCD动态扫描显示}

\begin{detailbox}[完整教学]
\textbf{【问题】}4位数码管需要4个BCD-7段译码器吗？如果每个译码器单独驱动→需要4×7=28个IO口。

\textbf{【方案】}动态扫描：1个译码器驱动所有数码管的段，快速轮流点亮各数码管（>50Hz），利用人眼视觉暂留。

\textbf{【电路】}
计数器→2-4译码器→轮流选通4个数码管的共阴/共阳极。BCD译码器输出段码到所有数码管。

\textbf{【VHDL核心逻辑】}
\begin{lstlisting}
process(clk)  -- 扫描时钟(通常>200Hz)
begin
  if rising_edge(clk) then
    scan_cnt <= scan_cnt + 1;
    case scan_cnt is
      when "00" => digit_sel <= "0001"; seg_data <= digit0;
      when "01" => digit_sel <= "0010"; seg_data <= digit1;
      when "10" => digit_sel <= "0100"; seg_data <= digit2;
      when "11" => digit_sel <= "1000"; seg_data <= digit3;
    end case;
  end if;
end process;
\end{lstlisting}

\textbf{【常见考题】}设计4位动态扫描显示电路。
\end{detailbox}


\section{变式2：带锁存的BCD译码器}

\begin{detailbox}[完整教学]
\textbf{【问题】}BCD译码器输入随时钟快速变化→数码管闪烁。需要"定格"显示。

\textbf{【方案】}在译码器前加寄存器。LE=1时锁存当前BCD值，LE=0时保持上次锁存值不变。

\textbf{【VHDL】}
\begin{lstlisting}
process(clk)
begin
  if rising_edge(clk) then
    if LE = '1' then
      latched_bcd <= bcd_input;  -- 锁存
    end if;
  end if;
end process;
seg <= bcd_to_7seg(latched_bcd);  -- 译码锁存后的值
\end{lstlisting}

\textbf{【设计思想】}锁存=寄存器=1个时钟沿保存数据。本质=在数据通路上插入一级寄存器。
\end{detailbox}


\section{变式3：共阳极 vs 共阴极数码管}

\begin{detailbox}[完整教学]
\textbf{【共阳极】}所有LED阳极接VCC。段驱动需低电平(0)点亮LED。seg='0'=亮。
\textbf{【共阴极】}所有LED阴极接GND。段驱动需高电平(1)点亮LED。seg='1'=亮。

\textbf{【输出值关系】}$seg_{\text{共阳}} = \text{not } seg_{\text{共阴}}$（互为反码）

\textbf{【VHDL适配】}
\begin{lstlisting}
-- 共阴极(1亮)
when "0000" => seg <= "1111110";  -- 0
-- 共阳极(0亮)
when "0000" => seg <= "0000001";  -- 0(反码)
\end{lstlisting}

\textbf{【考试陷阱】}题目未明确说明共阴/共阳时，看段码值判断：1111110(共阴0) vs 0000001(共阳0)。
\end{detailbox}


\section{变式4：计数器+译码器+数码管——综合题}

\begin{detailbox}[完整教学]
\textbf{【经典考题】}设计0-9计数器+BCD-7段译码器+数码管显示。

\textbf{【完整VHDL】}
\begin{lstlisting}
entity counter_display is
  port (clk, rst_n : in std_logic;
        seg : out std_logic_vector(6 downto 0));
end counter_display;

architecture rtl of counter_display is
  signal cnt : std_logic_vector(3 downto 0);
  function bcd7seg(bcd : std_logic_vector(3 downto 0))
    return std_logic_vector is
  begin
    case bcd is
      when "0000" => return "1111110";
      when "0001" => return "0110000";
      when "0010" => return "1101101";
      when "0011" => return "1111001";
      when "0100" => return "0110011";
      when "0101" => return "1011011";
      when "0110" => return "1011111";
      when "0111" => return "1110000";
      when "1000" => return "1111111";
      when "1001" => return "1111011";
      when others => return "0000000";
    end case;
  end function;
begin
  process(clk, rst_n)
  begin
    if rst_n='0' then cnt <= "0000";
    elsif rising_edge(clk) then
      if cnt = "1001" then cnt <= "0000";
      else cnt <= cnt + 1;
      end if;
    end if;
  end process;
  seg <= bcd7seg(cnt);
end rtl;
\end{lstlisting}

\textbf{【考点】}两个模块整合：计数器+译码器。
\end{detailbox}


\section{变式5：16进制显示译码器（0-F）}

\begin{detailbox}[完整教学]
\textbf{【问题】}BCD译码器只显示0-9。如何显示0-F（全4位二进制）？

\textbf{【方案】}扩展case覆盖1010-1111，定义A-F的段码。

\begin{lstlisting}
-- 扩展BCD译码器支持0-F
when "1010" => return "1110111";  -- A
when "1011" => return "0011111";  -- b
when "1100" => return "1001110";  -- C
when "1101" => return "0111101";  -- d
when "1110" => return "1001111";  -- E
when "1111" => return "1000111";  -- F
\end{lstlisting}

\textbf{【考试】}有时问A-F的段码值，需要记住典型表示法。
\end{detailbox}
